% chapters/design/data_model.tex
% Sección 4.2: Modelo de datos (~2-3 páginas)
% ============================================================================
%
% GUÍA: Presentar el modelo de datos con dos perspectivas complementarias.
%
%   --- DIAGRAMA 3: Modelo Entidad-Relación ---
%   Diagrama E-R de la base de datos PostgreSQL. Entidades principales:
%     - User (con campos de anonimización para RGPD)
%     - AcademicYear, Subject, Group
%     - Enrollment (alumno-grupo)
%     - Role, PermissionBlock, PermissionOverride
%     - Exam, ExamModel, Problem, OCRZone
%     - ExamInstance (con campo de estado: máquina de estados)
%     - Annotation (anotaciones multimedia)
%     - Grade, ReviewRequest
%     - AuditLog
%     - Notification
%
%   \begin{figure}{fig:er_diagram}{Modelo Entidad-Relación de la base de datos}
%     \image{}{}{er_diagram}
%   \end{figure}
%
%   TODO: Crear la imagen figures/er_diagram.png (o .pdf).
%         Herramientas: pgModeler, dbdiagram.io, draw.io, DBeaver.
%         NOTA: Si el diagrama E-R es muy grande, poner la versión simplificada
%         aquí y la versión completa en el Apéndice C.
%
%   --- DIAGRAMA 4: Diagrama de clases del dominio ---
%   Diagrama de clases UML centrado en los modelos Django más relevantes.
%   No necesita ser exhaustivo: destacar las relaciones, herencia y los
%   campos más importantes. Agrupar por subsistema funcional.
%
%   \begin{figure}{fig:class_diagram}{Diagrama de clases del modelo de dominio}
%     \image{}{}{class_diagram}
%   \end{figure}
%
%   TODO: Crear la imagen figures/class_diagram.png (o .pdf).
%         Herramientas: PlantUML, draw.io, pyreverse (auto-genera desde Django).
%
%   --- CONTENIDO TEXTUAL ---
%   1. Describir las entidades principales y sus relaciones.
%   2. Explicar las decisiones de diseño:
%      - Borrado lógico (soft delete) para RGPD (\ref{rnf:gdpr:deletion}).
%      - Versionado con campos para concurrencia optimista.
%      - Desnormalizaciones justificadas (si las hay).
%   3. Describir el modelo de permisos:
%      - Bloques funcionales con acciones.
%      - Cadena de evaluación: rol global → sobrecarga asignatura → individual.
%      - Referencia a \ref{rf:users} y \ref{rnf:maint:perms}.
%   4. Explicar cómo se almacenan los ficheros:
%      - PDFs originales en MinIO, referenciados por clave en BD.
%      - Anotaciones binarias (audio, imágenes) en MinIO.
%      - Referencia a \ref{rnf:perf:files}.
% ============================================================================

% TODO: Redactar la sección del modelo de datos y crear los diagramas.
